\documentclass{article}
\usepackage{amsmath}
\usepackage{geometry}

\newcommand{\PoissonDist}[1]{\mathrm{Poisson}\left(#1\right)}
\newcommand{\NBDist}[2]{\mathrm{NB}\left(\mathrm{mean}=#1,
\mathrm{dispersion}=#2\right)}
\newcommand{\GammaDist}[2]{\mathrm{Gamma}\left(\mathrm{shape}=#1,
\mathrm{rate}=#2\right)}
\newcommand{\MultinomialDist}[2]{\mathrm{Multinomial}\left(#1; #2\right)}
\newcommand{\incPMF}{p^{(\mathrm{inc})}}
\newcommand{\tostPMF}{p^{(\mathrm{tost})}}
\newcommand{\RAC}{\mathrm{RAC}}
\newcommand{\RAT}{\mathrm{RAT}}
\newcommand{\NBPMF}[3]{p_{\mathrm{NB}}\left(#1; \mathrm{mean}=#2,
\mathrm{dispersion}=#3\right)}
\newcommand{\Prob}[1]{\mathrm{Prob}\left(#1\right)}

\begin{document}

\section{Assumptions and notation}

Only symptom-onset counts \(C_t\) are observed. Infection counts \(I_t\)
and infectiousness variables \(Y_t,Z_t\), where present, are latent.
Naive models identify transmission and case events without explicitly
modelling incubation. Every infection eventually produces symptoms,
and there are no future importations.

\(\RAC(t)\) and \(\RAT(t)\) are the probabilities of at least one
additional case (symptom onset) and transmission, respectively, strictly
after day \(t\). Unless stated otherwise, risks condition on onsets through
\(t\) and fixed model parameters. Additional conditioning on latent histories
is stated locally; observation-only risks then require posterior averaging.
Distributional statements below specify sequential conditional models,
with conditioning on the relevant history suppressed.

The normalized delay PMFs are \(w_s\) (naive-model serial interval),
\(\incPMF_s\) (incubation period), and \(\tostPMF_s\) (time from
symptom onset to transmission), with supports \(1,\ldots,m\),
\(1,\ldots,m_{\mathrm{inc}}\), and \(0,\ldots,m_{\mathrm{tost}}\),
respectively. For integer \(a\geq0\), their CDFs are
\begin{equation}
  F_a^{(w)}=\sum_{s=1}^{a}w_s,\qquad
  F_a^{(\mathrm{inc})}=\sum_{s=1}^{a}\incPMF_s,\qquad
  F_a^{(\mathrm{tost})}=\sum_{s=0}^{a}\tostPMF_s,
\end{equation}
with PMFs zero outside their supports and empty sums equal to zero.
In particular, \(F_0^{(\mathrm{inc})}=0\) and
\(F_0^{(\mathrm{tost})}=\tostPMF_0\).
We use \(s\) for delays, \(a\) for one-based cohort age
(\(a=1\) for the cohort on day \(t\)), and \(\tau\) for calendar time.

\(\lambda_t\) denotes daily infectiousness before multiplication by
\(R_t\), and \(\Lambda_t^{(\mathrm{rem})}\) its accumulated remaining
contribution from existing cases. In SO models, \(\mu_t\)
is the conditional infection intensity and \(\bar N_t^{(\mathrm{inc})}\)
the conditional mean number still incubating at the end of day \(t\).
The bar denotes the expectation of \(N_t^{(\mathrm{inc})}\) under the
conditioning stated for the corresponding derivation.
Where assumed, constant future reproduction means
\(R_\tau=R^{(\mathrm{fut})}\) for every \(\tau>t\).

\(\NBDist{\mu}{\kappa}\) has variance \(\mu+\mu^2/\kappa\),
for \(\kappa>0\). Gamma distributions use shape and rate parameters.
Zero-shape Gamma variables and zero-intensity SSE negative-binomial
variables are defined as point masses at zero.

\section{Naive models}
\subsection{Poisson model}
\subsubsection{Model definition}

\begin{align}
  \lambda_{t} &= \sum_{s=1}^{m} C_{t-s} w_s,\\
  C_t &\sim \PoissonDist{R_{t}\lambda_{t}}.
\end{align}

\subsubsection{Risk of additional cases and transmission}

\begin{equation}
  \RAC(t) =\RAT(t)=
  1-\exp\left(-\sum_{\tau=t+1}^{t+m}R_{\tau}\sum_{s=\tau-t}^{m}
  C_{\tau-s} w_s\right).
\end{equation}

Assuming that \(R_\tau=R^{(\mathrm{fut})}\) for every \(\tau>t\), we have
\begin{align}
  \RAC(t)&=\RAT(t)=1-\exp(-R^{(\mathrm{fut})}\Lambda^{(\mathrm{rem})}_{t}),\\
  \Lambda^{(\mathrm{rem})}_{t} &= \sum_{a=1}^{m} (1-F_{a-1}^{(w)})C_{t+1-a}.
\end{align}

\subsection{DLO model}
\subsubsection{Model definition}

\begin{align}
  \lambda_{t} &= \sum_{s=1}^{m} C_{t-s} w_s,\\
  C_t &\sim \NBDist{R_{t}\lambda_{t}}{k}.
\end{align}

\subsubsection{Risk of additional cases and transmission}

\begin{align}
  \nonumber&\RAC(t)=\RAT(t)\\
  &=
  1-\prod_{\tau=t+1}^{t+m}\NBPMF{0}{R_{\tau}\sum_{s=\tau-t}^{m}
  C_{\tau-s} w_s}{k}\\
  &=
  1-\prod_{\tau=t+1}^{t+m}{\left(\frac{k}{k+R_{\tau}\sum_{s=\tau-t}^{m}
  C_{\tau-s} w_s}\right)}^{k}.
\end{align}

Assuming that \(R_\tau=R^{(\mathrm{fut})}\) for every \(\tau>t\), we have
\begin{align}
  \RAC(t)&=
  1-\prod_{\tau=t+1}^{t+m}{\left(\frac{k}{k+R^{(\mathrm{fut})}\sum_{s=\tau-t}^{m}
  C_{\tau-s} w_s}\right)}^{k}.
\end{align}

\subsection{SSE model}
\subsubsection{Model definition}

\begin{align}
  \lambda_{t} &= \sum_{s=1}^{m} C_{t-s} w_s,\\
  C_t &\sim \NBDist{R_{t}\lambda_{t}}{k\lambda_{t}}.
\end{align}

\subsubsection{Risk of additional cases and transmission}

\begin{align}
  \nonumber&\RAC(t)=\RAT(t)\\
  &=
  1-\prod_{\tau=t+1}^{t+m}\NBPMF{0}{R_{\tau}\sum_{s=\tau-t}^{m}
  C_{\tau-s} w_s}{k\sum_{s=\tau-t}^{m}
  C_{\tau-s} w_s}\\
  &=
  1-\prod_{\tau=t+1}^{t+m}{\left(\frac{k}{k+R_{\tau}}\right)}^{k\sum_{s=\tau-t}^{m}
  C_{\tau-s} w_s}.
\end{align}

Assuming that \(R_\tau=R^{(\mathrm{fut})}\) for every \(\tau>t\), we have
\begin{align}
  \RAC(t)&=
  1-{\left(\frac{k}{k+R^{(\mathrm{fut})}}\right)}^{k\Lambda^{(\mathrm{rem})}_{t}},\\
  \Lambda^{(\mathrm{rem})}_{t} &= \sum_{a=1}^{m} (1-F_{a-1}^{(w)})C_{t+1-a}.
\end{align}

\subsection{SSI model}
\subsubsection{Model definition}

\begin{align}
  \lambda_{t} &= \sum_{s=1}^{m} Y_{t-s} w_s,\\
  C_t &\sim \PoissonDist{R_{t}\lambda_{t}},\\
  Y_t &\sim \GammaDist{kC_t}{k}.
\end{align}

\subsubsection{Risk of additional cases and transmission}

Additionally conditioning on \(Y_\tau\) for \(\tau\leq t\):
\begin{equation}
  \RAC(t) =\RAT(t)=
  1-\exp\left(-\sum_{\tau=t+1}^{t+m}R_{\tau}\sum_{s=\tau-t}^{m}
  Y_{\tau-s} w_s\right).
\end{equation}

Assuming that \(R_\tau=R^{(\mathrm{fut})}\) for every \(\tau>t\), we have
\begin{align}
  \RAC(t)&=\RAT(t)=1-\exp(-R^{(\mathrm{fut})}\Lambda^{(\mathrm{rem})}_{t}),\\
  \Lambda^{(\mathrm{rem})}_{t} &= \sum_{a=1}^{m} (1-F_{a-1}^{(w)})Y_{t+1-a}.
\end{align}

\section{Symptom onset-anchored models}

On day \(t\), first obtain \(C_t\) from incubation allocations of earlier
infections, then sample any required infectiousness latents and \(I_t\),
using onsets through \(t\) and the TOST distribution. Finally, allocate
these infections to future onset days. Here \(I_{t,s}\) counts infections
on day \(t\) with incubation period \(s\):
\begin{align}
  C_t &= \sum_{s=1}^{m_{\mathrm{inc}}} I_{t-s,s},\\
  I_{t,1}, \ldots, I_{t,m_{\mathrm{inc}}} &\sim
  \MultinomialDist{I_t}{\incPMF_{1},
  \ldots, \incPMF_{m_{\mathrm{inc}}}}.
\end{align}
The equivalent forms below integrate out infection counts and incubation
allocations, retaining any infectiousness latents in the sequential model.
In each model, \(q_0\) denotes the probability that an individual with
onset after \(t\) produces no direct infections, integrating out its
future infectiousness.

\subsection{Poisson-SO model}
\subsubsection{Model definition}

\begin{align}
  \lambda_{t} &= \sum_{s=0}^{m_{\mathrm{tost}}} C_{t-s} \tostPMF_{s},\\
  I_t &\sim \PoissonDist{R_{t}\lambda_{t}}.
\end{align}

\subsubsection{Equivalent form}

\begin{align}
  C_t &\sim \PoissonDist{\sum_{s=1}^{m_{\mathrm{inc}}}
  \mu_{t-s} \incPMF_{s}},\\
  \mu_t &= R_{t}\sum_{s=0}^{m_{\mathrm{tost}}}
  C_{t-s} \tostPMF_{s}.
\end{align}

\subsubsection{Risk of additional cases and transmission}

Assume that \(R_\tau=R^{(\mathrm{fut})}\) for every \(\tau>t\).
Let \(N_t^{(\mathrm{rem})}\) count direct infections after \(t\) generated
by cases with onset by \(t\), and \(N_t^{(\mathrm{inc})}\) count individuals
infected by \(t\) whose onsets are after \(t\). Conditional on observed
onsets and parameters, these counts are independent, with
\begin{align}
  N^{(\mathrm{rem})}_{t} &\sim
  \PoissonDist{R^{(\mathrm{fut})}\Lambda^{(\mathrm{rem})}_{t}},\\
  \Lambda^{(\mathrm{rem})}_{t}&=\sum_{a=1}^{m_{\mathrm{tost}}}
  (1-F^{(\mathrm{tost})}_{a-1})C_{t+1-a},\\
  N^{(\mathrm{inc})}_{t} &\sim \PoissonDist{\bar N_t^{(\mathrm{inc})}},\\
  \bar N_t^{(\mathrm{inc})} &= \sum_{a=1}^{m_{\mathrm{inc}}}
  (1-F^{(\mathrm{inc})}_{a-1})\mu_{t+1-a}.
\end{align}
Poisson splitting leaves independent, unrevealed infection components
with onset after \(t\); nonnegative TOST prevents these individuals from
affecting the observed history before onset.

These counts exclude subsequent descendants, but their joint absence
is equivalent to no future cases. Hence
\begin{equation}
  \RAC(t)=1-\exp\left(-\left(R^{(\mathrm{fut})}\Lambda^{(\mathrm{rem})}_{t}+\bar N_t^{(\mathrm{inc})}\right)\right).
\end{equation}

Each incubating individual has its full transmission profile ahead of it,
with probability
\begin{equation}
  q_{0}=e^{-R^{(\mathrm{fut})}}
\end{equation}
of no direct transmissions.
Conditioning on the incubating count therefore gives
\begin{align}
  \nonumber&\RAT(t)\\
  &=
  1-\Prob{N^{(\mathrm{rem})}_{t}=0}\sum_{j=0}^{\infty}\Prob{\mathrm{no\ transm.\ from\ }j\ \mathrm{inc.\ individuals}}\Prob{N^{(\mathrm{inc})}_{t}=j}\\
  &=
  1-\exp\left(-R^{(\mathrm{fut})}\Lambda^{(\mathrm{rem})}_{t}\right)\sum_{j=0}^{\infty}q_{0}^{j}\Prob{N^{(\mathrm{inc})}_{t}=j}\\
  &=
  1-\exp\left(-\left(R^{(\mathrm{fut})}\Lambda^{(\mathrm{rem})}_{t}+\left(1-q_{0}\right)\bar N_t^{(\mathrm{inc})}\right)\right),
\end{align}
where the sum can be evaluated either directly or by noting that this
is the probability generating function of a Poisson distribution evaluated
at \(q_{0}\).

\subsection{SSE-SO model}
\subsubsection{Model definition}

\begin{align}
  \lambda_{t} &= \sum_{s=0}^{m_{\mathrm{tost}}} C_{t-s} \tostPMF_{s},\\
  I_t &\sim \NBDist{R_{t}\lambda_{t}}{k\lambda_{t}}.
\end{align}

\subsubsection{Equivalent form}

First rewrite in terms of a Poisson-Gamma mixture:
\begin{align}
  \lambda_{t} &= \sum_{s=0}^{m_{\mathrm{tost}}} C_{t-s} \tostPMF_{s},\\
  Z_t &\sim \GammaDist{k\lambda_{t}}{k},\\
  I_t &\sim \PoissonDist{R_{t}Z_t}.
\end{align}

Equivalent to:
\begin{align}
  C_t &\sim \PoissonDist{\sum_{s=1}^{m_{\mathrm{inc}}}
  \mu_{t-s} \incPMF_{s}},\\
  \mu_t &= R_{t}Z_t,\\
  \lambda_{t} &= \sum_{s=0}^{m_{\mathrm{tost}}} C_{t-s} \tostPMF_{s},\\
  Z_t &\sim \GammaDist{k\lambda_{t}}{k}.
\end{align}

\subsubsection{Risk of additional cases and transmission}

Assume that \(R_\tau=R^{(\mathrm{fut})}\) for every \(\tau>t\).
Additionally conditioning on \(Z_\tau\) for \(\tau\leq t\), the counts
\(N_t^{(\mathrm{rem})}\) and \(N_t^{(\mathrm{inc})}\) are independent, with
\begin{align}
  N^{(\mathrm{rem})}_{t} &\sim
  \NBDist{R^{(\mathrm{fut})}\Lambda^{(\mathrm{rem})}_{t}}{k\Lambda^{(\mathrm{rem})}_{t}},\\
  \Lambda^{(\mathrm{rem})}_{t}&=\sum_{a=1}^{m_{\mathrm{tost}}}
  (1-F^{(\mathrm{tost})}_{a-1})C_{t+1-a},\\
  N^{(\mathrm{inc})}_{t} &\sim \PoissonDist{\bar N_t^{(\mathrm{inc})}},\\
  \bar N_t^{(\mathrm{inc})} &= \sum_{a=1}^{m_{\mathrm{inc}}}
  (1-F^{(\mathrm{inc})}_{a-1})\mu_{t+1-a}.
\end{align}

Then, by similar reasoning to the Poisson-SO model, we have
\begin{align}
  \RAC(t)&=1-\left(\frac{k}{k+R^{(\mathrm{fut})}}\right)^{k\Lambda_{t}^{(\mathrm{rem})}}\exp\left(-\bar N_t^{(\mathrm{inc})}\right),\\
  \RAT(t)&=1-\left(\frac{k}{k+R^{(\mathrm{fut})}}\right)^{k\Lambda_{t}^{(\mathrm{rem})}}\exp\left(-\left(1-q_{0}\right)\bar N_t^{(\mathrm{inc})}\right),
\end{align}
where here, since the offspring distribution for future cases is
\(\NBDist{R^{(\mathrm{fut})}}{k}\), we have
\begin{align}
  q_{0}&=\NBPMF{0}{R^{(\mathrm{fut})}}{k} =
  \left(\frac{k}{k+R^{(\mathrm{fut})}}\right)^{k}.
\end{align}

\subsection{SSI-SO model}
\subsubsection{Model definition}

\begin{align}
  \lambda_{t} &= \sum_{s=0}^{m_{\mathrm{tost}}} Y_{t-s} \tostPMF_s,\\
  I_t &\sim \PoissonDist{R_{t}\lambda_{t}},\\
  Y_t &\sim \GammaDist{kC_t}{k}.
\end{align}

\subsubsection{Equivalent form}

\begin{align}
  C_t &\sim \PoissonDist{\sum_{s=1}^{m_{\mathrm{inc}}}
  \mu_{t-s} \incPMF_{s}},\\
  \mu_t &= R_{t}\sum_{s=0}^{m_{\mathrm{tost}}}
  Y_{t-s} \tostPMF_{s},\\
  Y_t &\sim \GammaDist{kC_t}{k}.
\end{align}

\subsubsection{Risk of additional cases and transmission}

Assume that \(R_\tau=R^{(\mathrm{fut})}\) for every \(\tau>t\).
Additionally conditioning on \(Y_\tau\) for \(\tau\leq t\), the counts
\(N_t^{(\mathrm{rem})}\) and \(N_t^{(\mathrm{inc})}\) are independent, with
\begin{align}
  N^{(\mathrm{rem})}_{t} &\sim
  \PoissonDist{R^{(\mathrm{fut})}\Lambda^{(\mathrm{rem})}_{t}},\\
  \Lambda^{(\mathrm{rem})}_{t}&=\sum_{a=1}^{m_{\mathrm{tost}}}
  (1-F^{(\mathrm{tost})}_{a-1})Y_{t+1-a},\\
  N^{(\mathrm{inc})}_{t} &\sim \PoissonDist{\bar N_t^{(\mathrm{inc})}},\\
  \bar N_t^{(\mathrm{inc})} &= \sum_{a=1}^{m_{\mathrm{inc}}}
  (1-F^{(\mathrm{inc})}_{a-1})\mu_{t+1-a}.
\end{align}

Then,
\begin{align}
  \RAC(t)&=1-\exp\left(-\left(R^{(\mathrm{fut})}\Lambda^{(\mathrm{rem})}_{t}+\bar N_t^{(\mathrm{inc})}\right)\right),\\
  \RAT(t)&=1-\exp\left(-\left(R^{(\mathrm{fut})}\Lambda^{(\mathrm{rem})}_{t}+\left(1-q_{0}\right)\bar N_t^{(\mathrm{inc})}\right)\right),
\end{align}
where, integrating out the Gamma infectiousness of each future case,
\begin{equation}
  q_{0}=\NBPMF{0}{R^{(\mathrm{fut})}}{k}
  =\left(\frac{k}{k+R^{(\mathrm{fut})}}\right)^k.
\end{equation}

\end{document}
